% ============================================================
%  04_technologies.tex  –  Chapter 3: Technologies Used
% ============================================================

\chapter{Technologies Used}
\label{ch:technologies}

\subsection*{Raspberry Pi 4 Model B}

Raspberry Pi 4 Model B is a single-board computer developed by the Raspberry Pi
Foundation~\cite{rpi4specs}. In this project it has the following uses: central processing
unit, reading sensor data, computing gas concentrations, applying correction factors and
warning generation.

\vspace{12pt}
\begin{figure}[H]
  \centering
  \includegraphics[width=0.4\textwidth]{rpi4}
  \caption{Figure 1: Raspberry Pi 4 Model B (\url{https://www.robot-advance.com/EN/art-raspberry-pi-4-model-b-4go-2640.htm})}
  \label{fig:rpi4}
\end{figure}
\vspace{12pt}

\subsection*{Sensors}

\textbf{DHT11} is a low-cost digital sensor for measuring temperature and humidity. It
communicates with Raspberry Pi through a single-wire digital interface and is directly
connected to a GPIO pin (pin 16). A 5\,k$\Omega$ pull-up resistor is placed between the VCC
and the data lines to guarantee signal stability. It works on 3.3\,V and its readings are used
in monitoring and correcting MQ sensor PPM calculations.

\vspace{12pt}
\begin{figure}[H]
  \centering
  \includegraphics[width=0.3\textwidth]{dht11}
  \caption{Figure 2: DHT11 sensor (\url{https://www.optimusdigital.ro/ro/senzori-senzori-de-temperatura/99-senzor-de-temperatura-si-si-umiditate-dht11.html})}
  \label{fig:dht11}
\end{figure}
\vspace{12pt}

\textbf{MQ-135} is an analog gas sensor sensitive to a broad range of air pollutants, used in
this project primarily as an indicator of overall air quality and estimated CO\textsubscript{2}
equivalent (eCO\textsubscript{2}) concentration. It works on 5\,V and the output is an analog
voltage in proportion to gas concentration. Raspberry Pi lacks native analog input, so the
sensor is connected through the MCP3208 by a voltage divider. The divider is formed of a
1\,k$\Omega$ and 2\,k$\Omega$ resistor combination. All voltage dividers from this project are
of the same build.

\vspace{12pt}
\begin{figure}[H]
  \centering
  \includegraphics[width=0.3\textwidth]{mq135}
  \caption{Figure 3: MQ-135 sensor (\url{https://sigmanortec.ro/Senzor-calitate-aer-MQ-135-p126101726})}
  \label{fig:mq135}
\end{figure}
\vspace{12pt}

The \textbf{MQ-3} sensor is analog. Its function is to detect alcohol vapor and benzene. The
sensor is interfaced to the MCP3208 ADC using a resistor voltage divider. The sensor works at
5\,volts.

\vspace{12pt}
\begin{figure}[H]
  \centering
  \includegraphics[width=0.3\textwidth]{mq3}
  \caption{Figure 4: MQ-3 sensor (\url{https://sigmanortec.ro/Senzor-gaz-MQ-3-Alcool-Etanol-p126101225})}
  \label{fig:mq3}
\end{figure}
\vspace{12pt}

\textbf{MQ-6}, an analog sensor, is used to detect LPG and butane gases. It helps monitor
gas leaks in a home environment. Connected through MCP3208 ADC by using the resistor voltage
divider and it works on 5\,V.

\vspace{12pt}
\begin{figure}[H]
  \centering
  \includegraphics[width=0.3\textwidth]{mq6}
  \caption{Figure 5: MQ-6 sensor (\url{https://www.optimusdigital.ro/en/gas-sensors/1132-modul-senzor-de-gaz-mq-6.html})}
  \label{fig:mq6}
\end{figure}
\vspace{12pt}

\textbf{MQ-7}, an analog sensor, detects carbon monoxide (CO), a colorless and odorless gas
that poses a serious health risk even at low concentrations. Connected through the MCP3208 ADC
by using a resistor voltage divider and it operates on 5\,V.

\vspace{12pt}
\begin{figure}[H]
  \centering
  \includegraphics[width=0.3\textwidth]{mq7}
  \caption{Figure 6: MQ-7 sensor (\url{https://uk.robotshop.com/products/mq-7-gas-sensor})}
  \label{fig:mq7}
\end{figure}
\vspace{12pt}

\textbf{MQ-8} is an analog sensor which detects hydrogen gas (H\textsubscript{2}). Connected
through the MCP3208 ADC by using a resistor voltage divider and it operates on 5\,V.

\vspace{12pt}
\begin{figure}[H]
  \centering
  \includegraphics[width=0.3\textwidth]{mq8}
  \caption{Figure 7: MQ-8 sensor (\url{https://www.optimusdigital.ro/en/gas-sensors/1131-modul-senzor-de-gaz-mq-8.html})}
  \label{fig:mq8}
\end{figure}
\vspace{12pt}

\textbf{KY-037} is a sound detection module that provides both a digital output (DOUT) and an
analog output (AOUT). The analog output is read through the MCP3208 to measure the background
ambient sound level as a voltage, while the digital output is connected directly to a GPIO pin
and monitored via hardware interrupt to count discrete sound events within each 5-second
measurement cycle.

\vspace{12pt}
\begin{figure}[H]
  \centering
  \includegraphics[width=0.3\textwidth]{ky037}
  \caption{Figure 8: KY-037 sensor (\url{https://ardushop.ro/en/modules/1589-high-sensitivity-sound-detection-microphone-module-ky-037-6427854023735.html})}
  \label{fig:ky037}
\end{figure}
\vspace{12pt}

\subsection*{MCP3208}

The MCP3208 is an 8-channel, 12-bit analog-to-digital converter (ADC) that communicates with
the Raspberry Pi over the SPI (Serial Peripheral Interface) protocol~\cite{mcp3208}. Since the
Raspberry Pi GPIO pins only accept digital signals, the MCP3208 is essential to make the
connection between analog MQ sensors and the KY-037 sound sensor and the Raspberry Pi. It works
on 3.3\,V and is connected to Raspberry Pi's SPI pins (CLK on pin 23, MOSI on pin 19, MISO on
pin 21, and CS on pin 24).

\vspace{12pt}
\begin{figure}[H]
  \centering
  \includegraphics[width=0.35\textwidth]{mcp3208}
  \caption{Figure 9: MCP3208 (\url{https://www.digikey.ro/en/products/detail/microchip-technology/MCP3208-CI-P/305928})}
  \label{fig:mcp3208}
\end{figure}
\vspace{12pt}

\subsection*{Additional Hardware Components}

Two LEDs (blue and red) are used as simple visual indicators. The blue LED signals when data is
sent to the server. The red LED lights up when a warning condition is detected.

A button is used to start or stop the sensor readings when pressed.

\subsection*{Python}

Python is a programming language which focuses on code readability, simplicity, ease of
writing, while also offering a large library. It also supports a wide range of applications,
from web development to data analysis and machine learning.

\subsection*{Flask}

Flask is a lightweight Python web framework. It was used to build the backend of the
project~\cite{flask}. It is designed to make scalability of apps easy.

\subsection*{PostgreSQL and AWS RDS}

PostgreSQL is an open-source relational database system~\cite{postgresql}. In this project it
stores all sensor readings, registered user accounts, form messages and site settings. The
database is hosted on Amazon Web Services (AWS) Relational Database Service
(RDS)~\cite{awsrds}. This manages the database related part without requiring manual server
administration.

\subsection*{AWS EC2}

Amazon Web Services Elastic Compute Cloud (EC2) is a cloud computing service~\cite{awsec2}. It
provides virtual server instances accessible over the internet. In this project, an EC2 instance
hosted in the EU (Stockholm) region runs the Flask server continuously, making the website and
the sensor data API permanently available online. The EC2 instance serves both the static
website files and handles all API requests from the Raspberry Pi and from website users.

\subsection*{HTML}

Hypertext Markup Language (HTML) is the standard markup language for documents designed to be
displayed in a web browser. It is used to define the content and structure of web content.

\subsection*{CSS}

Cascading Style Sheets (CSS) is a stylesheet language used to describe the presentation of a
document written in HTML or XML. CSS is used to render elements on screen, paper, speech or on
other media.

\subsection*{Tailwind CSS}

Tailwind CSS is a utility-first CSS framework that offers pre-built classes for styling HTML
elements directly~\cite{tailwind}. In this project, Tailwind CSS was not used as the primary
styling framework, however, specific UI components were adopted from Tailwind CSS for a faster
implementation.

\subsection*{JavaScript}

JavaScript (JS) is a lightweight interpreted (or just-in-time compiled) programming language.
It is mostly used as the scripting language for Web pages, but many non-browser environments
also use it, such as Node.js, Apache CouchDB and Adobe Acrobat.

\subsection*{Google Colab}

Google Colab is a free cloud-based Jupyter notebook environment provided by Google, offering
access to CPU and GPU compute resources without requiring local
installation~\cite{googlecolab}. In this project, Google Colab was used as the training
environment for both machine learning models.

\subsection*{Flutter}

Flutter is an open-source UI toolkit created by Google, used for developing applications for
mobile, web and desktop. The application code is written in Dart and compiled into native ARM
binary code for Android. In this project, Flutter was used for creating the \airnest{} Android
app that replicates the functionality of certain pages of the website.
